\subsubsection{三卷分工}

\begin{center}
\begin{tabular}{@{}>{\raggedright\arraybackslash}p{0.21\textwidth}>{\raggedright\arraybackslash}p{0.43\textwidth}>{\raggedright\arraybackslash}p{0.24\textwidth}@{}}
\toprule[1.5pt]
卷 & 内容 & 何时读 \\
\midrule[1pt]
卷一 · 基础通识 & Linux、C++、CMake、Git & 入门必过 \\
卷二 · 专项技术栈 & ROS2、通信、SDR、OpenCV、深度学习、CUDA、相机、滤波、BA、MPC、导航 & 按方向选修 \\
卷三 · 本卷 & 赛场链路、接口边界、工程准则 & 开营精读 + 中期答辩 \\
\bottomrule[1.5pt]
\end{tabular}
\end{center}

\begin{example}[把问题放回正确卷册]
建立三行书签：卷一编译与 CMake、卷二几何视觉与 BA、卷三工程准则。分别给“链接错误”“PnP 点序”“联调前如何留证据”连线。检查：先按问题性质找入口，避免从头翻整套书。
\end{example}

\begin{exercise}
\label{v3-ex-index-1}
Git 操作不熟、坐标广播积压、调车没有版本记录，分别优先查哪卷？
\end{exercise}

\subsubsection{按方向的推荐顺序}

\textbf{自瞄（学习时间 $\rightarrow$ 工程时间）：}
卷一 $\rightarrow$ 本卷第\ref{ch:v3-overview}、\ref{ch:v3-aim}、\ref{ch:v3-engineering}章 $\rightarrow$ 卷二 ROS2、OpenCV、深度学习、相机、滤波、BA、MPC $\rightarrow$ 回本卷做链路答辩 $\rightarrow$ 实车仓库。

\textbf{雷达：}
卷一 $\rightarrow$ 本卷第\ref{ch:v3-overview}、\ref{ch:v3-radar}、\ref{ch:v3-engineering}章 $\rightarrow$ 卷二 ROS2、多机通信、SDR、OpenCV、深度学习、相机 $\rightarrow$ 开源/队内雷达仓库部署文档。

\textbf{导航：}
卷一 $\rightarrow$ 本卷第\ref{ch:v3-overview}、\ref{ch:v3-nav}、\ref{ch:v3-engineering}章 $\rightarrow$ 卷二 ROS2、导航（通信建议）$\rightarrow$ 建图与实车。

\begin{example}[给导航新人排顺序]
写出“卷一通识 → 本卷职责、导航、工程准则 → 卷二 ROS2 与导航 → 回到本卷答辩 → 受控实车”。在每一步旁标注输出：基础小工程、链路图、专项练习、现场演示。检查：学习计划包含产出，不能只有书单。
\end{example}

\begin{exercise}
\label{v3-ex-index-2}
雷达新人不会串口帧解析，是否应先因“看起来高级”而跳去 MPC？写出更直接的补课顺序。
\end{exercise}

\subsubsection{资料索引（活页）}

\begin{itemize}
    \item 卷二各章：\texttt{02-专项技术栈/}
    \item 自瞄公式与导读：\texttt{视觉培训/最小二乘法/}（讲义 + BA/Kalman/MPC）
    \item 雷达站参考实现：队内指定开源仓库与部署说明（注意许可证与权重策略）
    \item 相机 / 控制器 PPT、视觉小课堂：队内资料目录
    \item 裁判系统协议：当年官方文档（勿依赖过期拷贝）
\end{itemize}

\begin{example}[检查一个真实入口]
在本地文件浏览器依次打开项目目录、02-专项技术栈、main.tex，搜索“多机通信方法”，沿其 input 找到章节文件。记录入口、章名和检查日期。若队内讲义路径在本机不存在，记录“本机未取得访问入口”，交负责人确认。检查：书中提到的路径不等于当前电脑必定存在。
\end{example}

\begin{note}
队内仓库地址、接口责任人和复盘日期由负责人填写。本书不虚构这些信息；本机没有路径时应记录待确认，不要把未访问的资料标成已核验。
\end{note}

\begin{conclusion}
\begin{itemize}
  \item 路径回答在哪里，版本回答适用哪套实现。
  \item 责任人回答由谁确认，核对日期回答何时验证。
  \item 四项一起记录，才能让索引可追溯。
\end{itemize}
\end{conclusion}

\begin{exercise}
\label{v3-ex-index-3}
只有一个没有版本号的协议截图，够不够作为联调依据？还需取得什么信息？
\end{exercise}

\subsubsection{版本与责任}

\begin{center}
\begin{tabular}{@{}p{0.30\textwidth}p{0.62\textwidth}@{}}
\toprule[1.5pt]
项目 & 填写 \\
\midrule[1pt]
教材版本 & \texttt{2026.08} \\
本卷维护 & 视觉组培训负责人 \\
接口文档责任人 & （队内填写） \\
下次复盘日期 & （队内填写） \\
\bottomrule[1.5pt]
\end{tabular}
\end{center}
本章为活页：仅入口路径变化时更新本章即可；协议字段或行为变化时，还须同步接口文档、实现和验证，不能只换链接。

\begin{example}[完成一次索引变更]
教学情景：部署说明从路径 A 迁到 B。记录旧入口、新入口、适用分支或版本、检查日期与维护人；用新入口实际打开文档，再让另一位同学复核。无法访问 B 时记录受阻原因，不能标成已完成。
\end{example}

\begin{exercise}
\label{v3-ex-index-4}
教材活页改了协议链接，但消息字段定义也变了，仅修改活页是否足够？
\end{exercise}
